\section{引言}
\label{sec:intro}
\subsection{问题与目标}
给定有向图 $G=(V,E)$、非负边权 $w:E\to\mathbb{Z}_{\ge0}$ 和源点 $s$，
求从 $s$ 到各顶点的最短距离 $d[v]$。不可达顶点的距离记作 $\infty$。
记 $n=|V|$、$m=|E|$；本文的 $m$ 始终指有向弧条数。
课程项目要求使用至少两种堆结构且包含 Fibonacci 堆，
在 USA 路网中至少评测 1000 对查询，并给出运行时间、规模曲线与分析。

Dijkstra 的贪心顺序依赖最小优先队列。二叉堆、段树及 Fibonacci 堆
对插入、减键和取最小元素的处理方式不同，因而具有不同的渐近代价和实际常数。
对非负整数距离，还可以使用单调的 radix heap。
本报告关注两个问题：理论上更便宜的减键操作能否转化为实际收益；
固定图结构和输入后，不同实现的速度次序如何变化。

\subsection{实现与实验范围}
七种 Dijkstra 队列为标准库优先队列、手写二叉堆、段树堆、Fibonacci 堆、
radix heap、GNU PBDS pairing heap 和 thin heap。
除标准库版本以重复插入模拟减键外，其他版本支持按顶点定位并更新键。
FIFO SPFA 与 SLF+LLL SPFA 是基于队列的标签修正算法，不属于 Dijkstra 的堆变体，
在报告中作为独立参照。

所有实现共用 \code{general-dijkstra.hpp} 中的算法主循环以及 CSR 图。
随机图与固定密度实验对同组参数使用固定种子 42 和同一个生成规则。
路网采用 \code{USA-road-d.NY} 与 \code{USA-road-d.BAY}，
查询文件是项目生成的 1000 对随机查询，种子为 971；查询文件不是官方原始下载。
本轮未测试 NE、E 或全美 USA 图，结论不外推到这些规模。

\subsection{本轮结果的解释边界}
尽管路网输入格式为点对点查询，现有算法会从每个源点计算完整距离数组，
直到队列为空，目标点仅用于取校验和。
因此本文的“ms/查询”是\textbf{一个点对触发一次完整单源搜索的平均时间}，
不能等同于遇到目标点即停止的点对点搜索延迟。
本轮不包含路径恢复、预处理索引或多线程并行搜索。

随机图和复杂度实验每组重复三轮；每个路网实现仅运行一个包含 1000 查询的批次。
所有图表和表格都来自同一轮已校验的 CSV 快照，避免把旧报告中的测量值混入比较。
